\DocumentMetadata{
  pdfversion=2.0,
  pdfstandard=ua-2,
  lang=en-US,
  tagging=on,
  tagging-setup={math/setup=mathml-SE,tabsorder=structure}
}
\documentclass[11pt,letterpaper]{article}
\usepackage[margin=0.68in,includefoot]{geometry}
\usepackage{newcomputermodern}
\usepackage{amsmath,amscd,mathtools}
\usepackage{tikz,adjustbox}
\providecommand{\square}{\mdlgwhtsquare}
\usepackage[hidelinks]{hyperref}
\hypersetup{
  pdftitle={Algebra PhD Exam, May 1989},
  pdfauthor={University of Florida Department of Mathematics},
  pdfsubject={Graduate mathematics examination},
  pdfdisplaydoctitle=true,
  bookmarksopen=true
}
\setcounter{secnumdepth}{0}
\setlength{\parindent}{0pt}
\setlength{\parskip}{0.28em}
\setlength{\emergencystretch}{3em}
\setlength{\leftmargini}{2.15em}
\setlength{\leftmarginii}{2.65em}
\setlength{\leftmarginiii}{3em}
\pagestyle{empty}
\raggedbottom
\allowdisplaybreaks
\NewTaggingSocketPlug{block/list/label}{examproblem}{%
  \tagstructbegin{tag=Lbl}\tagstructend%
  \tagstructbegin{tag=\LBody}%
  \tagstructbegin{tag=\ProblemHeadingTag,title-o={\ProblemHeadingText}}%
  \tagmcbegin{tag=\ProblemHeadingTag,actualtext={\ProblemHeadingText}}%
  #2%
  \tagmcend%
  \tagstructend%
}
\newenvironment{examproblems}[2]{%
  \def\ProblemHeadingTag{#2}%
  \AssignTaggingSocketPlug{block/list/label}{examproblem}%
  \begin{enumerate}%
  \setcounter{enumi}{#1}%
  \renewcommand{\labelenumi}{\textbf{\theenumi.}}%
  \setlength{\topsep}{0.35em}%
  \setlength{\partopsep}{0pt}%
  \setlength{\itemsep}{0.48em}%
  \setlength{\parsep}{0.18em}%
}{\end{enumerate}}
\newenvironment{examsubparts}{%
  \AssignTaggingSocketPlug{block/list/label}{default}%
  \begin{enumerate}%
  \setlength{\topsep}{0.18em}%
  \setlength{\partopsep}{0pt}%
  \setlength{\itemsep}{0.16em}%
  \setlength{\parsep}{0.1em}%
}{\end{enumerate}}
\newenvironment{examinstructions}{%
  \par\begingroup\itshape
}{%
  \par\endgroup\vspace{0.18em}
}
\makeatletter
\renewcommand\subsection{\@startsection{subsection}{2}{0pt}%
  {-1.25ex plus -.25ex minus -.1ex}%
  {0.55ex plus .1ex}%
  {\normalfont\large\bfseries}}
\renewcommand{\@maketitle}{%
  \newpage\null\vskip -1em%
  {\footnotesize\bfseries
    \href{https://www.ufl.edu/}{University of Florida}, \href{https://math.ufl.edu/}{Department of Mathematics}\par}%
  \vskip -0.5em%
  \tagstructbegin{tag=H1,title={Algebra PhD Exam, May 1989}}%
  \tagpdfparaOff%
  \tagmcbegin{tag=H1}%
    {\LARGE\bfseries\href{https://gma.math.ufl.edu/exams/algebra-phd/}{Algebra PhD Exam}, May 1989\par}%
  \tagmcend%
  \tagpdfparaOn%
  \tagstructend%
  \vskip 0.25em%
  \tagmcbegin{artifact}%
    \hrule height 1pt%
  \tagmcend%
  \vskip 1em%
}
\makeatother
\title{Algebra PhD Exam, May 1989}
\author{}
\date{}
\begin{document}
\maketitle
\thispagestyle{empty}
\begin{examinstructions}
Do 7 out of the following questions.
\end{examinstructions}
\begin{examproblems}{0}{H2}
\def\ProblemHeadingText{Problem 1.}
\item
This problem has four parts.
\begin{examsubparts}
\item[\textnormal{(a)}]
Define: Solvable group.
\item[\textnormal{(b)}]
State the Burnside’s theorem for solvability of a finite group.
\item[\textnormal{(c)}]
Let~\(H\) be a normal subgroup of a group~\(G\). Prove that \(C_G(H)\) is a normal subgroup of~\(G\).
\item[\textnormal{(d)}]
Let~\(G\) be a finite group of order 12. Prove that~\(G\) has a non-trivial normal 2-subgroup (i.e., there is a subgroup~\(T\) of~\(G\) such that \(1\ne T\), and~\(T\) is a 2-subgroup, and~\(T\) is normal in~\(G\)).
\end{examsubparts}
\def\ProblemHeadingText{Problem 2.}
\item
Two permutation representations \(\phi_i:G\to\Sigma(X_i)\), (the symmetric group on \(X_i\)), are equivalent if there exists \(\alpha:X_1\to X_2\) such that for all \(g\in G\), \(x\in X_1\), we have \(x(g\phi_1)\alpha=x(\alpha)(g\phi_2)\). Prove or disprove that the following two permutation representations of \(S_4\) of degree 6 are equivalent:
\begin{examsubparts}
\item[\textnormal{(a)}]
on the right coset space of \(\langle(12)(34),(13)(24)\rangle\).
\item[\textnormal{(b)}]
on the right coset space of \(\langle(12),(34)\rangle\).
\end{examsubparts}
\def\ProblemHeadingText{Problem 3.}
\item
Let \(F_n\) be the free group on~\(n\) generators \(x_1,\ldots,x_n\). Let \(F_n'\) denote the commutator subgroup. Prove that \(F_n/F_n'\) is the free abelian group on \(\{x_1,\ldots,x_n\}\).
\def\ProblemHeadingText{Problem 4.}
\item
In each of the following groups decide whether it is a semi-direct product of two of its proper subgroups: \(D_4\) (the dihedral group of order 8); \(Q_8\) (the quaternion group of order 8). Justify your answer.
\def\ProblemHeadingText{Problem 5.}
\item
Let~\(A\) be a commutative ring with identity. Suppose that for each prime ideal~\(P\), the local ring \(A_P\) has no nilpotent element \(\ne 0\). Show that~\(A\) has no nilpotent element \(\ne 0\). If each \(A_P\) is an integral domain, is~\(A\) necessarily an integral domain? Justify your answer.
\def\ProblemHeadingText{Problem 6.}
\item
Let~\(A\) be a commutative ring with identity, \(a\) an ideal, \(M\) an~\(A\)-module.
\begin{examsubparts}
\item[\textnormal{(a)}]
Prove that \(A\otimes_A M\cong M\) as~\(A\)-modules.
\item[\textnormal{(b)}]
Prove that \((A/a)\otimes_A M\) is isomorphic to \(M/aM\).
\end{examsubparts}
\def\ProblemHeadingText{Problem 7.}
\item
Let \(A\subseteq B\) be integral domains, \(B\) integral over~\(A\). Prove that~\(B\) is a field if and only if~\(A\) is a field.
\def\ProblemHeadingText{Problem 8.}
\item
Find a minimal primary decomposition for the following ideals. Also give the prime ideals to which the primary ideals belong.
\begin{examsubparts}
\item[\textnormal{(a)}]
\((36)\) as an ideal in~\(\mathbb{Z}\).
\item[\textnormal{(b)}]
\((x^2,xy)\) as an ideal in \(\mathbb{C}[x,y]\).
\end{examsubparts}
\def\ProblemHeadingText{Problem 9.}
\item
State the following theorems.
\begin{examsubparts}
\item[\textnormal{(a)}]
Nakayama’s lemma.
\item[\textnormal{(b)}]
Hilbert basis theorem.
\item[\textnormal{(c)}]
Jacobson’s Density theorem.
\item[\textnormal{(d)}]
Maschke’s theorem.
\item[\textnormal{(e)}]
Wedderburn’s theorem for finite dimensional simple algebras (or Wedderburn–Artin’s theorem for simple Artinian rings).
\end{examsubparts}
\end{examproblems}
\end{document}
